% Authoritative exercise source. Preserve all derivations and diagrams.

\ExerciseStart{G3.08}{Seguire CNOT su base e sovrapposizione}{Consolidamento}

\BookHeading{Problema, dati e obiettivo}

Applicare \(\mathrm{CNOT}_{0\to1}\) alla preparazione \((3\ket{00}-\ket{10}+2i\ket{11})/\sqrt{14}\). Verificare prima la porta sull’intera base, quindi calcolare tutte le probabilità e le marginali.

\ExTheory{\ExRef{c3-registro}{3.07}}{La porta agisce come \(\ket{a,b}\mapsto\ket{a,b\oplus a}\). Conserva il controllo e cambia il bersaglio soltanto nel sottospazio con \(a=1\). La linearità mantiene le ampiezze coerenti.}

\ExPrediction{Le componenti 10 e 11 devono scambiarsi. La probabilità del controllo non dovrebbe cambiare, mentre la distribuzione del bersaglio può cambiare.}

\ExFigure{\begin{tikzpicture}[>=Stealth,exlabel]
\node[exbox] (n0) at (0.0,0.0) {$00\mapsto00$};\node[exbox] (n1) at (3.6,0.0) {$01\mapsto01$};\draw[exother] (n0)--(n1);\node[exbox] (n2) at (7.2,0.0) {$10\leftrightarrow11$};\draw[exother] (n1)--(n2);
\end{tikzpicture}}{Permutazione dei quattro vettori di base.}


\begin{ExTasks}

\item Costruire la tabella della porta.

\item Scrivere la matrice nella base dichiarata.

\item Applicare ogni riga al vettore iniziale.

\item Calcolare probabilità e marginali.

\item Applicare nuovamente la porta e controllare il ritorno.

\end{ExTasks}

\ExSolution

\ExStep{Verificare la base}

\ExEq{00&\mapsto00,\quad01\mapsto01,\\10&\mapsto11,\quad11\mapsto10.}
Il controllo non viene misurato durante l’operazione.

\ExStep{Costruire la matrice}

\ExEq{C&=\begin{pmatrix}1&0&0&0\\0&1&0&0\\0&0&0&1\\0&0&1&0\end{pmatrix}.}
Le colonne sono le immagini dei quattro vettori di base.

\ExStep{Applicare tutte le righe}

\ExEq{C\ket\psi&=\frac1{\sqrt{14}}\colv{1\cdot3\\1\cdot0\\1\cdot2i\\1(-1)}
=\frac1{\sqrt{14}}\colv{3\\0\\2i\\-1}.}
La norma iniziale e quella finale valgono entrambe \((9+0+4+1)/14=1\).

\ExStep{Calcolare la lettura}

\ExEq{P(00)&=9/14,\quad P(01)=0,\\P(10)&=|2i|^2/14=4/14,\quad P(11)=1/14.}
Le marginali del controllo sono \(9/14\) e \(5/14\), come prima. Quelle del bersaglio diventano \(13/14\) e \(1/14\), rispetto ai valori iniziali \(10/14\) e \(4/14\).

\ExStep{Verificare l’inversa}

La seconda CNOT scambia di nuovo le ultime due componenti e restituisce \((3,0,-1,2i)^T/\sqrt{14}\). Ogni vettore di base torna a se stesso, quindi \(C^2=I\) su tutto il registro.

\par\noindent\begin{minipage}{\linewidth}
\ExFigure{\begin{tikzpicture}[>=Stealth,exlabel]
\draw[exaxis] (-.6,0)--(4.4,0) node[right] {esito};\draw[exaxis] (-.6,0)--(-.6,3.6) node[above] {$p$};\node[left] at (-.6,3) {$ 1 $};\draw[Ink!15,dashed] (-.6,3)--(3.9000000000000004,3);\filldraw[fill=Accent!25,draw=Accent] (0.0,0) rectangle (0.65,1.9285714285714288);\node[below] at (0.325,0) {$ 00 $};\node[above] at (0.325,1.9285714285714288) {$ 0.6429 $};\filldraw[fill=Accent!25,draw=Accent] (1.1,0) rectangle (1.75,0.0);\node[below] at (1.425,0) {$ 01 $};\node[above] at (1.425,0.0) {$ 0 $};\filldraw[fill=Accent!25,draw=Accent] (2.2,0) rectangle (2.85,0.8571428571428571);\node[below] at (2.5250000000000004,0) {$ 10 $};\node[above] at (2.5250000000000004,0.8571428571428571) {$ 0.2857 $};\filldraw[fill=Accent!25,draw=Accent] (3.3000000000000003,0) rectangle (3.95,0.21428571428571427);\node[below] at (3.6250000000000004,0) {$ 11 $};\node[above] at (3.6250000000000004,0.21428571428571427) {$ 0.07143 $};
\end{tikzpicture}}{Probabilità dopo CNOT; le etichette conservano l’ordine big-endian.}


\ExInterpret{Il controllo coerente applica una trasformazione lineare allo stato congiunto. La conservazione della marginale del controllo e della norma offre due controlli distinti dell’implementazione.}

\ExLab{Nella sezione 3.07 scegliere «CNOT su (3,0,-1,2i)/\(\sqrt{14}\)». Confrontare le quattro ampiezze prima e dopo la porta; verificare esplicitamente la fase della componente trasferita da 11 a 10.}


\end{minipage}
